\documentclass[12pt,a4paper]{article}

\usepackage[margin=1in]{geometry}
\usepackage{amsmath}
\usepackage{graphicx}
\usepackage{xcolor}
\usepackage{booktabs}
\usepackage{hyperref}

\title{A Short Technical Report}
\author{Technical Writing Practical Example}
\date{\today}

\begin{document}

\maketitle
\tableofcontents
\newpage

\section{Introduction}
This report demonstrates basic LaTeX document preparation, including
sections, mathematical expressions, tables, figures and cross-references.

\section{Mathematical Typesetting}
The quadratic equation is
\begin{equation}
  ax^2 + bx + c = 0.
  \label{eq:quadratic}
\end{equation}
Its solutions are
\[
  x = \frac{-b \pm \sqrt{b^2 - 4ac}}{2a}.
\]

A matrix can be typeset as
\[
A = \begin{bmatrix}
1 & 2 \\
3 & 4
\end{bmatrix}.
\]

\section{Tables}
Table~\ref{tab:results} presents sample values.

\begin{table}[htbp]
\centering
\caption{Sample experimental values}
\label{tab:results}
\begin{tabular}{lrr}
\toprule
Trial & Input & Output \\
\midrule
1 & 2 & 4 \\
2 & 3 & 9 \\
3 & 4 & 16 \\
\bottomrule
\end{tabular}
\end{table}

\section{Figures}
Figure~\ref{fig:placeholder} demonstrates a caption and label.
\begin{figure}[htbp]
\centering
\fbox{\parbox[c][3cm][c]{0.65\linewidth}{\centering Placeholder for a technical illustration}}
\caption{A framed placeholder figure}
\label{fig:placeholder}
\end{figure}

\section{Cross-References}
Equation~\ref{eq:quadratic} gives the quadratic form.
The data are listed in Table~\ref{tab:results}, and the illustration
appears in Figure~\ref{fig:placeholder}.

\section{Conclusion}
This small document demonstrates common LaTeX structures used in
technical reports.

\end{document}
